\documentclass[10pt,a4paper]{amsart}
\usepackage[margin=19mm]{geometry}
\usepackage{amsmath,amssymb,mathtools}
\usepackage[scr=rsfs,cal=euler]{mathalfa}
\usepackage{xfrac}
\usepackage[unicode,hidelinks,bookmarks=false]{hyperref}
\newcommand{\ie}{i.e.\ }
\newcommand{\cf}{cf.\ }
\newcommand{\resp}{respectively\ }
\newcommand{\Subsection}{\subsection}
\providecommand{\nobreakdash}{\nobreak\hbox{-}}
\setlength{\emergencystretch}{2em}
\multlinegap0pt
\usepackage{bm}
\usepackage{bbm}
\usepackage{dsfont}
\usepackage{xspace}
\usepackage{cancel}
\usepackage{mathtools}
\usepackage{amsthm}
\makeatletter
\@ifundefined{thm}{\newtheorem{thm}{Thm}}{}
\@ifundefined{claim}{\newtheorem{claim}[thm]{Claim}}{}
\@ifundefined{prop}{\newtheorem{prop}[thm]{Proposition}}{}
\makeatother
\newcommand{\N}{\mathbb{N}}
\newcommand{\Prob}{\operatorname{Prob}}
\newcommand{\Sub}{\operatorname{Sub}}
\newcommand{\rL}{\operatorname{L}}
\title[Strong primeness for equivalence relations arising from Zari]{Strong primeness for equivalence relations arising from Zariski dense subgroups}
\author{Daniel Drimbe, Cyril Houdayer}
\date{Source: arXiv:2407.01806v2 (CC BY 4.0)}
\begin{document}
\maketitle

\noindent\emph{Source and licence.} Strong primeness for equivalence relations arising from Zariski dense subgroups. Version arXiv:2407.01806v2 (CC BY 4.0), distributed under the Creative Commons Attribution 4.0 International licence, as recorded in the arXiv licence metadata for this exact version and shown on the abstract page https://arxiv.org/abs/2407.01806v2. Full rights notice: licenses/arxiv-2407.01806-CC-BY-4.0.txt.\par\smallskip

\noindent\textit{Editorial scope.} Counted unit: the Prop whose source label is \texttt{prop-borel} in the article (source0.tex lines 709--713, source label \texttt{prop-borel}), delivered together with the article's own complete original proof of it (source0.tex lines 715--748, one \texttt{proof} environment of 34 lines). No mathematics is written, repaired or re-derived: statement and proof are the article's own text line for line. Required context retained uncounted: none. Every definition, display and label of those lines is retained unchanged. Omitted from this package and not redistributed: 1--708, 714--714, 749--1118. Nothing inside a retained excerpt is omitted or altered: every retained body is delivered exactly as the article's lines are written, so no omission receipt of any kind accompanies this package. Citations: the retained excerpts carry no citation command at all, so no bibliography entry of the article is transcribed and no reference is redistributed. Reference adaptation: none was needed and none was made; every reference inside the delivered unit resolves inside it, so no unresolved reference or citation is produced. Printed slips kept literal: no printed word, symbol, hypothesis or number of the counted statement or of its proof was altered. Layout and class: the article's own class is \texttt{amsart} with packages including amsmath, amsthm, amsfonts, amssymb, mathrsfs, mathtools, enumerate, hyperref, graphicx, caption, subcaption, pgf, tikz, appendix; the wrapper is the lane's pinned \texttt{amsart} editorial wrapper at 10pt on A4, which restates the theorem-like environments the retained text needs and adds the article's own macro definitions that the retained text needs (\texttt{\textbackslash N}, \texttt{\textbackslash Prob}, \texttt{\textbackslash Sub}, \texttt{\textbackslash rL}), with extra wrapper packages (bm, bbm, dsfont, xspace, cancel, mathtools). The delivered numbering is the unit's own. Assets: no retained span contains a figure environment, a graphics-inclusion command, a PDF-page inclusion, a table or a TikZ picture; no image file is copied into this package, the package declares no asset dependency and no image file is redistributed with it. Licence: version arXiv:2407.01806v2 of this article is distributed under the Creative Commons Attribution 4.0 International licence, as recorded in the arXiv licence metadata for this exact version and shown on the abstract page https://arxiv.org/abs/2407.01806v2. A full rights notice is delivered at licenses/arxiv-2407.01806-CC-BY-4.0.txt. Characters: every retained excerpt is pure ASCII, so the delivered engine, XeLaTeX with its default Latin Modern text font, reports no missing character.
\begin{prop}\label{prop-borel}
Keep the same notation as above. Then the map 
$$\Psi : \Sub(G) \to \Sub(\mathscr G) : H \mapsto \rL^0(X, \nu, H)$$ 
is Borel.
\end{prop}

\begin{proof}
Consider the Wijsman topology $\tau = \tau_{W(d_{\mathscr G})}$ on $\Sub(\mathscr G)$. Fix a countable dense subset $D \subset \mathscr G$. For every $f \in D$ and all $\alpha, \beta \in \mathbb Q \cap [0, 1]$ such that $\alpha < \beta$, define the open subsets 
\begin{align*}
O(f, \alpha)^+ &= \left \{\mathscr K \in \Sub(\mathscr G) \mid d_{\mathscr G}(f, \mathscr K) > \alpha \right \} \\
O(f, \beta)^- &= \left \{\mathscr K \in \Sub(\mathscr G) \mid d_{\mathscr G}(f, \mathscr K) < \beta \right \} \\
O(f, \alpha, \beta) &= O(f, \alpha)^+ \cap O(f, \beta)^-.
\end{align*}
Then the countable set $\left \{ O(f, \alpha, \beta) \mid f \in D, \alpha, \beta \in \mathbb Q \cap [0, 1],\alpha < \beta \right \}$ is a subbase for the topology $\tau$ on $\Sub(\mathscr G)$. Thus, in order to show that the map $\Psi : \Sub(G) \to \Sub(\mathscr G)$ is Borel, it suffices to show that $\Psi^{-1}(O(f, \alpha)^+)$ and $\Psi^{-1}(O(f, \alpha)^-)$ are Borel subsets of $\Sub(G)$ for every $f \in D$ and every $\alpha \in \mathbb Q \cap [0, 1]$.


\begin{claim}\label{claim-lusin}
Let $(H_n)_{n \in \N}$ be a sequence in $\Sub(G)$ and $H \in \Sub(G)$ such that $H_n \to H$ in $\Sub(G)$. Set $\mathscr H_n = \rL^0(X, \nu, H_n)$ for every $n \in \N$ and $\mathscr H = \rL^0(X, \nu, H)$. Then for every $f \in \mathscr H$, there exists a sequence $(f_n)_{n \in \N}$ in $\mathscr G$ such that $f_n \in \mathscr H_n$ for every $n \in \N$ and $\lim_n d_{\mathscr G}(f, f_n) = 0$.
\end{claim}

Indeed, let $f \in \mathscr H$. Without loss of generality, we may assume that $X$ is a compact metrizable space and $\nu \in \Prob(X)$ is a Borel probability measure. Regard $f : X \to H$ as a measurable function. By Lusin's theorem, for every $n \geq 1$, there exists a closed subset $X_n \subset X$ such that $\nu(X_n) \geq 1 - 1/n$ and $f|_{X_n} : X_n \to H$ is continuous. By compactness, there exist $r \geq 1$ and $h_1, \dots, h_r \in H$ such that $X_n \subset \bigcup_{j = 1}^r f^{-1}(B_{d}(h_j, 1/n))$, where $B_d(h, \varepsilon) \subset G$ denotes the open ball in $G$ of center $h \in G$ and radius $\varepsilon > 0$. Set $Y_1 = X_n \cap f^{-1}(B_{d}(h_1, 1/n))$ and $Y_j = X_n \cap f^{-1}(B_{d}(h_j, 1/n)) \setminus \bigcup_{i = 1}^{j - 1} Y_i$ for every $2 \leq j \leq r$. Then $(Y_j)_{1 \leq j \leq r}$ is a measurable partition of $X_n$. Since $H_n \to H$ in $\Sub(G)$, for every $1 \leq j \leq r$, we may find $h_j^n \in B_{d}(h_j, 1/n) \cap H_n$. Define the measurable map $f_n : X \to H_n$ by $f_n|_{X \setminus X_n} = e \mathbf 1_{X \setminus X_n}$ and $f_n|_{Y_j} = h_j^n \mathbf 1_{Y_j}$ for every $1 \leq j \leq r$. Then we have
\begin{align*}
\limsup_n d_{\mathscr G}(f, f_n) &= \limsup_n \int_X d(f(x), f_n(x)) \,{\rm d}\nu(x) \\
&\leq \limsup_n \left(\nu(X \setminus X_n) + \frac2n \nu(X_n) \right) =0.
\end{align*}
This finishes the proof of Claim \ref{claim-lusin}.

\begin{claim}\label{claim-closed}
Let $(H_n)_{n \in \N}$ be a sequence in $\Sub(G)$ and $H \in \Sub(G)$ such that $H_n \to H$ in $\Sub(G)$. Set $\mathscr H_n = \rL^0(X, \nu, H_n)$ for every $n \in \N$ and $\mathscr H = \rL^0(X, \nu, H)$. Then for every $f \in \mathscr G$, we have
$$\limsup_n d_{\mathscr G}(f, \mathscr H_n) \leq d_{\mathscr G}(f, \mathscr H).$$
\end{claim}

Indeed, set $\alpha = d_{\mathscr G}(f, \mathscr H)$ and $\beta = \limsup_n d_{\mathscr G}(f, \mathscr H_n)$. Choose an increasing sequence $(n_j)_{j \in \N}$ in $\N^*$ such that $\beta = \lim_j d_{\mathscr G}(f, \mathscr H_{n_j})$. Choose a sequence $(f_n)_{n \in \N}$ in $\mathscr H$ such that $\alpha \leq d_{\mathscr G}(f, f_n) < \alpha + 1/n$ for every $n \geq 1$. By Claim \ref{claim-lusin}, we may choose a sequence $(h_n)_{n \geq 1}$ in $\mathscr G$ such that $h_n \in \mathscr H_n$ and $d_{\mathscr G}(f_n, h_n) \leq 1/n$ for every $n \geq 1$. In particular, for every $j \in \N$, we have 
$$d_{\mathscr G}(f, \mathscr H_{n_j}) \leq d_{\mathscr G}(f, h_{n_j}) \leq d_{\mathscr G}(f, f_{n_j}) + d_{\mathscr G}(f_{n_j}, h_{n_j}) \leq \alpha + \frac{2}{n_j}.$$
This implies that $\beta \leq \alpha$ and finishes the proof of Claim \ref{claim-closed}.

Let $f \in D$ and $\alpha \in \mathbb Q \cap [0, 1]$. Claim \ref{claim-closed} implies that $\Psi^{-1}(O(f, \alpha)^-) \subset \Sub(G)$ is open. Observe that 
$$O(f, \alpha)^+ = \bigcup_{j = 1}^\infty  \left \{\mathscr K \in \Sub(\mathscr G) \mid d_{\mathscr G}(f, \mathscr K) \geq \alpha + 1/j \right \}.$$
Claim \ref{claim-closed} implies that $\Psi^{-1}(O(f, \alpha)^+) \subset \Sub(G)$ is a countable union of closed subsets and so $\Psi^{-1}(O(f, \alpha)^+) \subset \Sub(G)$ is Borel. Therefore, we have showed that the map $\Psi : \Sub(G) \to \Sub(\mathscr G) $ is Borel.
\end{proof}

\end{document}
